\documentclass[11pt,a4paper]{article}
\usepackage[margin=27mm,headheight=15pt]{geometry}
\usepackage{fontspec}
\usepackage{xeCJK}
\setCJKmainfont{FandolSong-Regular.otf}[BoldFont=FandolSong-Bold.otf]
\usepackage{amsmath,amssymb,amsthm,mathtools}
\usepackage{booktabs,array,tabularx}
\usepackage{xcolor}
\usepackage{enumitem}
\usepackage{fancyhdr}
\usepackage{hyperref}
\hypersetup{colorlinks=true,linkcolor=blue!45!black,
  urlcolor=blue!45!black,citecolor=blue!45!black,
  pdftitle={A counterexample to the explicit multiplicity-gap criterion for Conjecture 00000001421},
  pdfsubject={Ordinary spectral gap, total-variation cutoff, and a Lean 4 certificate}}
\setlength{\parindent}{1.2em}
\setlength{\parskip}{0.18em}
\setlength{\emergencystretch}{2em}
\setlist{itemsep=0.3em,topsep=0.4em}
\pagestyle{fancy}
\fancyhf{}
\fancyhead[L]{\small Conjecture 00000001421}
\fancyhead[R]{\small Ordinary gap and total variation}
\fancyfoot[C]{\thepage}
\renewcommand{\headrulewidth}{0.3pt}
\newtheorem{theorem}{Theorem}[section]
\newtheorem{proposition}[theorem]{Proposition}
\newtheorem{lemma}[theorem]{Lemma}
\newtheorem{corollary}[theorem]{Corollary}
\theoremstyle{definition}
\newtheorem{definition}[theorem]{Definition}
\theoremstyle{remark}
\newtheorem*{remark}{Remark}
\newcommand{\TV}{\mathrm{TV}}
\newcommand{\C}{\mathbb C}
\newcommand{\R}{\mathbb R}
\newcommand{\N}{\mathbb N}
\newcommand{\Z}{\mathbb Z}
\newcommand{\ind}{\mathbf 1}
\newcommand{\tmix}{t_{\mathrm{mix}}}
\newcommand{\spec}{\operatorname{spec}}
\newcommand{\mult}{\operatorname{mult}}
\newcommand{\code}[1]{\texttt{\detokenize{#1}}}
\newcommand{\module}[1]{\texttt{\detokenize{#1}}}

\title{\vspace{-1.4em}\textbf{A counterexample to the explicit\\ multiplicity--gap criterion}\\[0.4em]
  \large Conjecture 00000001421}
\author{}
\date{}

\begin{document}
\maketitle
\vspace{-2.2em}

\begin{abstract}
We construct a sequence of exactly-half-lazy, uniform Cayley random walks on
finite abelian groups of increasing order. Its ordinary spectral gap is
$1/N$. The smallest nontrivial transition eigenvalue has multiplicity $N-1$;
the smallest positive eigenvalue of $I-P$ has multiplicity one. Consequently
both possible multiplicity-to-gap ratios diverge. Nevertheless the actual
total-variation distance at every positive integer time is
$\tfrac12(1-1/N)^t$. The true quarter-mixing time diverges, while its ratio to
the quarter-to-three-quarter transition width tends to one. The family has
no standard total-variation cutoff. This refutes the conjecture's explicit
numerical sufficient condition under the ordinary-gap and total-variation
interpretations. The source's undefined conditional-gap, second-order, and
comparison-language clauses are not assigned additional meanings. A closed
Lean~4 certificate proves the stochastic, spectral, and mixing-time bridges
used in the argument.
\end{abstract}

\section{Original statement and precise scope}

The original statement is reproduced in both supplied languages.

\begin{quote}\small
\textbf{English.} Definition: Cutoff means that for a family of random walks
the ratio of total mixing time to mixing window tends to infinity.
Conjecture: Lazy random walks on abelian groups exhibit cutoff if and only
if the second-order behavior of the conditional spectral gap over
irreducible characters of the generating support satisfies an explicit
criterion: the ratio of the multiplicity of the minimal nontrivial
eigenvalue to the spectral gap tends to infinity; the criterion is
completely characterized by Diaconis--Saloff-Coste type second-eigenvalue
comparison.
\end{quote}

\begin{quote}\small
\textbf{中文。} 定义：cutoff 指随机游走族的总混合时间与混合窗之比趋于无穷。猜想：交换群上惰性随机游走出现 cutoff 当且仅当其生成支撑的不可约特征条件谱隙的二阶行为满足显式判据:最小非平凡特征值的重数与谱隙之比趋于无穷;判据由 Diaconis–Saloff-Coste 型第二特征值比较完全刻画。
\end{quote}

The explicit ratio criterion has a definite sufficient direction once the
gap, eigenvalue convention, and mixing metric are specified. We refute that
direction with the following scope.

\begin{enumerate}[label=\arabic*.]
\item \textbf{Gap and metric.} The gap is the ordinary spectral gap of the
reversible transition operator. Mixing is measured in total variation from
the uniform stationary law, with the maximum over starting states. The
source does not itself name a metric or give a precise window definition.
We state our definitions below and prove failure of both a genuine
fixed-threshold width test and the quantified standard TV cutoff criterion.
\item \textbf{Two eigenvalue readings.} We compute separately the
multiplicity of the smallest nontrivial eigenvalue of $P$ and that of the
smallest positive eigenvalue of $I-P$. The first is not silently replaced
by the second. Both resulting ratios diverge in the same counterexample.
\item \textbf{Characters and support.} The relevant characters are those
of the group generated by the support; the support itself need not be a
group. Our support generates the entire state group. Its genuine complex
group characters are the one-dimensional irreducible characters of this
finite abelian group, and they form the complete spectral basis used here.
\item \textbf{Undefined prose.} The statement does not define a
conditioning operation or a quantity called the ``conditional spectral
gap,'' a ``second-order behavior,'' or a particular Diaconis--Saloff-Coste
comparison assertion. We do not invent such definitions or claim a theorem
about an unspecified comparison principle. If an additional nonstandard
quantity or further restriction was intended, a clarified conjecture would
be required. The result here concerns the explicit numerical sufficient
criterion under the ordinary-gap and TV interpretations just stated.
\end{enumerate}

The counterexample satisfies the narrower, familiar convention of exactly
$1/2$ holding and otherwise uniform choice from a symmetric nonidentity
generating set. It uses neither variable holding probabilities nor a
weighted generating set. Its state-space orders and ordinary mixing times
both tend to infinity.

\section{Definitions and the counterexample}

Let $n\in\N$ and set $N=n+4$. Thus $N\ge4$ throughout, and every limit below
is taken along this sequence as $n\to\infty$, equivalently $N\to\infty$.
Write
\[
G_N=(\Z/2\Z)\times(\Z/N\Z),\qquad
S_N=\{(0,h):h\ne0\}\cup\{(1,0)\}.
\]
Define a discrete-time random walk by adding independent increments with law
\begin{equation}\label{eq:jump}
\mu_N(g)=\frac12\ind_{\{g=0\}}
        +\frac1{2N}\ind_{\{g\in S_N\}}.
\end{equation}
Its transition matrix and the candidate stationary law are
\begin{equation}\label{eq:kernel}
P_N(x,y)=\mu_N(y-x),\qquad \pi_N(y)=\frac1{2N}.
\end{equation}
Powers $P_N^t$ always mean the usual powers of this matrix, with $P_N^0=I$.

\begin{definition}\label{def:mixing}
For $t\in\N$ define
\[
d_N(t)=\max_{x\in G_N}\frac12\sum_{y\in G_N}
          \left|P_N^t(x,y)-\pi_N(y)\right|.
\]
For $\varepsilon>0$, the mixing time is the genuine least threshold time
\[
t_N(\varepsilon)=\min\{t\in\N:d_N(t)\le\varepsilon\},
\]
once nonemptiness of this set has been proved. The ordinary mixing time and
the quarter-to-three-quarter width are
\[
T_N=t_N(1/4),\qquad W_N=t_N(1/4)-t_N(3/4).
\]
The subtraction and ratios involving these times are taken in $\R$.
The standard quantified total-variation cutoff criterion used here is
\begin{equation}\label{eq:cutoff}
\text{for every }0<\varepsilon<\frac12,\qquad
\frac{t_N(\varepsilon)}{t_N(1-\varepsilon)}\longrightarrow1.
\end{equation}
For each fixed $\varepsilon>0$, the denominator is positive once
$N>1/(2\varepsilon)$, since $d_N(0)=1-1/(2N)>1-\varepsilon$.
Thus these asymptotic ratios are eventually defined; a convention at finitely
many initial indices does not affect their limits. Lean uses total real division.
\end{definition}

The fixed-tolerance ratio convention in \eqref{eq:cutoff} is standard; see
Levin's introductory definition, slide~23 of~\cite{levin}. That reference is
cited only to identify the convention. The construction and proof below
were derived independently before consulting it, and no proof content is
taken from it. We use non-strict threshold inequalities in the definition
of $t_N$, consistently with the Lean development.

\begin{theorem}\label{thm:main}
The matrices in \eqref{eq:kernel} give a growing family of irreducible,
reversible, exactly-half-lazy uniform Cayley walks. Put
\begin{equation}\label{eq:qrs}
q=1-\frac1N,\qquad r=\frac12,\qquad s=\frac12-\frac1N.
\end{equation}
Their complete transition spectrum, with algebraic multiplicity, is
\[
1\ (1),\qquad q\ (1),\qquad r\ (N-1),\qquad s\ (N-1).
\]
The ordinary spectral gap is $\delta_N=1/N$. The literal minimum
nontrivial-transition multiplicity and the minimum positive $I-P_N$
multiplicity satisfy, respectively,
\begin{equation}\label{eq:ratios}
\frac{N-1}{\delta_N}=N(N-1)\longrightarrow\infty,
\qquad
\frac1{\delta_N}=N\longrightarrow\infty.
\end{equation}
Nevertheless
\begin{equation}\label{eq:tv-main}
d_N(0)=1-\frac1{2N},\qquad
d_N(t)=\frac12q^t\quad(t\ge1).
\end{equation}
All positive-tolerance mixing times are finite and attained, and
\begin{equation}\label{eq:time-main}
t_N(3/4)=1,\qquad T_N\ge\frac N2,\qquad
W_N=T_N-1>0,\qquad \frac{T_N}{W_N}\longrightarrow1.
\end{equation}
In particular $T_N\to\infty$, the ratio $T_N/W_N$ does not tend to infinity,
and the family does not satisfy \eqref{eq:cutoff}.
\end{theorem}

The remaining sections prove every assertion of this theorem directly.

\section{The support and the stochastic walk}

The horizontal part of $S_N$ has $N-1$ elements and is disjoint from the
single vertical element $(1,0)$. Hence $|S_N|=N$ and $0\notin S_N$.
Negation preserves the horizontal part, while $-(1,0)=(1,0)$; thus
$S_N=-S_N$. Every $(a,b)\in G_N$ is
\[
(a,b)=(a,0)+(0,b).
\]
Each summand is either zero or belongs to $S_N$, because $a$ is either zero
or one. Consequently the subgroup generated by $S_N$ is all of $G_N$.
The group order is $|G_N|=2N=2(n+4)\to\infty$.

More explicitly, from $(a,b)$ the chain has probabilities
\[
\begin{array}{c|c}
\text{destination}&\text{probability}\\ \hline
(a,b)&1/2\\
(a,d),\ d\ne b&1/(2N)\ \text{for each }d\\
(a+1,b)&1/(2N).
\end{array}
\]
All other transition probabilities are zero. The probabilities are
nonnegative and each row sums to $1/2+N/(2N)=1$. Symmetry of the support and
increment law gives $P_N(x,y)=P_N(y,x)$, so the column sums are also one.
The law $\pi_N$ is nonnegative, has total mass one, and satisfies
\[
\sum_x\pi_N(x)P_N(x,y)=\frac1{2N}\sum_xP_N(x,y)=\pi_N(y).
\]
This symmetry also proves detailed balance with respect to $\pi_N$.

For any $x=(a,b)$ and $y=(c,d)$, choose $z=(a,d)$. Then
$P_N(x,z)>0$, since $x$ and $z$ have the same first coordinate, and
$P_N(z,y)>0$, since $z$ and $y$ have the same second coordinate and their
first coordinates are either equal or differ by one. Therefore
\[
P_N^2(x,y)=\sum_zP_N(x,z)P_N(z,y)>0.
\]
The chain is irreducible. Its diagonal entries are exactly $1/2$, so it is
lazy and aperiodic. Finally, products of nonnegative stochastic matrices
are stochastic: summing a row of $AB$ gives
$\sum_z A(x,z)\sum_yB(z,y)=1$. Thus every actual matrix power $P_N^t$ is a
probability kernel, as required in Definition~\ref{def:mixing}.

\section{Complete character spectra and genuine multiplicities}

\subsection{Characters form the entire eigenbasis}

A complex character of a finite additive group is a map $\chi$ satisfying
$\chi(0)=1$ and $\chi(u+v)=\chi(u)\chi(v)$. Its values have modulus one:
if $ku=0$, then $\chi(u)^k=1$.
The $N$ characters of $\Z/N\Z$ are
$\chi_j(b)=\exp(2\pi i jb/N)$, $0\le j<N$; indeed a character is determined
by its value at one, which must be an $N$th root of unity. Similarly the two
characters of $\Z/2\Z$ are $a\mapsto(-1)^{\epsilon a}$ for
$\epsilon\in\{0,1\}$.

Every character of $G_N$ factors uniquely as
\begin{equation}\label{eq:characters}
\psi_{\epsilon,\chi}(a,b)=(-1)^{\epsilon a}\chi(b).
\end{equation}
To see completeness of this list as characters, restrict any character
$\psi$ to the two factors and use
$\psi(a,b)=\psi(a,0)\psi(0,b)$. Uniqueness follows from the same restrictions.

For a nontrivial character $\eta$ of any finite group $H$, choose $h_0$
with $\eta(h_0)\ne1$. Translation of the sum gives
\[
\sum_{h\in H}\eta(h)=\sum_{h\in H}\eta(h+h_0)
=\eta(h_0)\sum_{h\in H}\eta(h),
\]
so the sum is zero. For the trivial character it is $|H|$.
For distinct characters, their quotient is nontrivial; this identity gives
orthogonality under the usual complex inner product on functions on $H$.
The $2N$ distinct characters in \eqref{eq:characters} are therefore
orthogonal nonzero vectors in the $2N$-dimensional space $\C^{G_N}$.
They form a basis. In particular, no spectral modes are omitted.

\subsection{Action of the actual transition matrix}

For a character $\psi$ and a state $x$, the matrix from \eqref{eq:kernel}
satisfies
\begin{align*}
(P_N\psi)(x)
&=\sum_y\mu_N(y-x)\psi(y)\\
&=\sum_g\mu_N(g)\psi(x+g)
 =\psi(x)\sum_g\mu_N(g)\psi(g).
\end{align*}
Thus the actual eigenvalue of $\psi_{\epsilon,\chi}$ is
\begin{equation}\label{eq:character-eigenvalue}
\lambda_{\epsilon,\chi}
=\frac12+\frac1{2N}\left(\sum_{h\ne0}\chi(h)+(-1)^\epsilon\right).
\end{equation}
The preceding sum identity and $\chi(0)=1$ give
\[
\sum_{h\ne0}\chi(h)=
\begin{cases}N-1,&\chi\text{ trivial},\\-1,&\chi\text{ nontrivial}.
\end{cases}
\]
Substitution into \eqref{eq:character-eigenvalue} yields
\begin{center}
\begin{tabular}{@{}llll@{}}
\toprule
First-factor character&Second-factor character&Eigenvalue&Multiplicity\\
\midrule
Trivial&Trivial&$1$&$1$\\
Nontrivial&Trivial&$q=1-1/N$&$1$\\
Trivial&Nontrivial&$r=1/2$&$N-1$\\
Nontrivial&Nontrivial&$s=1/2-1/N$&$N-1$\\
\bottomrule
\end{tabular}
\end{center}
Because this is a complete basis of eigenvectors, these counts are the
algebraic multiplicities of the actual operator. Equivalently its
characteristic polynomial over $\C$ is
\begin{equation}\label{eq:charpoly}
\det(XI-P_N)=(X-1)(X-q)(X-r)^{N-1}(X-s)^{N-1}.
\end{equation}
For $N\ge4$ the values are strictly ordered:
\begin{equation}\label{eq:ordering}
0<s<r<q<1.
\end{equation}
Consequently the only stationary character is the trivial one, and the
ordinary spectral gap is
\[
\delta_N=1-\max\{\lambda:\lambda\in\spec(P_N),\ \lambda\ne1\}
=1-q=\frac1N.
\]
All eigenvalues are positive, so using the largest absolute nontrivial
eigenvalue produces the same gap in this family.

\subsection{The distinct $I-P_N$ interpretation}

The same complete character basis diagonalizes the actual complementary
operator $L_N=I-P_N$, with eigenvalue $1-\lambda_{\epsilon,\chi}$ on each
character. Its characteristic polynomial is
\begin{equation}\label{eq:complement-charpoly}
\det(XI-L_N)
=X\bigl(X-(1-q)\bigr)
       \bigl(X-(1-r)\bigr)^{N-1}
       \bigl(X-(1-s)\bigr)^{N-1}.
\end{equation}
Thus its complete spectrum, with multiplicity, is
\[
0\ (1),\qquad \frac1N\ (1),\qquad
\frac12\ (N-1),\qquad \frac12+\frac1N\ (N-1).
\]
The smallest positive eigenvalue is precisely $\delta_N=1/N$, and its
algebraic multiplicity is one. In contrast, the smallest nontrivial
\emph{transition} eigenvalue is $s$, with multiplicity $N-1$.
Distinctness of the roots in \eqref{eq:charpoly} and
\eqref{eq:complement-charpoly} makes these exact root multiplicities,
rather than merely lower bounds supplied by some eigenvectors.
Both claims in \eqref{eq:ratios} now follow. Their divergence follows from
$N\to\infty$, since $N(N-1)\ge N$ for $N\ge4$.

\section{An identity for the actual matrix powers}

We derive the transition probabilities directly by matrix multiplication.
This also supplies a proof of the bridge between the stochastic walk and
the scalar expressions used for its mixing analysis.

For $x=(a,b)$ and $y=(c,d)$ put
\[
U=\ind_{\{a=c\}},\qquad V=\ind_{\{b=d\}},\qquad
\sigma=2U-1\in\{1,-1\}.
\]
For four real numbers $A,B,C,D$, let
\[
F(A,B,C,D)(x,y)=A+BU+CV+D\ind_{\{x=y\}}.
\]
This is an auxiliary matrix, not a new definition of the walk. Write
$c_N=1/(2N)$. The original transition kernel has the decomposition
\begin{equation}\label{eq:decomp}
P_N(x,y)=s\ind_{\{x=y\}}+c_N U+c_N V.
\end{equation}
When $x=y$ the right side is $s+2c_N=1/2$; in the two possible nonidentity
neighbor cases it is $c_N$; otherwise it is zero. This verifies
\eqref{eq:decomp} directly from \eqref{eq:kernel}.

The coordinate sums of $F$ are
\begin{align}
\sum_{d\in\Z/N\Z}F(A,B,C,D)(x,(c,d))
  &=NA+C+(NB+D)\ind_{\{a=c\}},\label{eq:sum-second}\\
\sum_{c\in\Z/2\Z}F(A,B,C,D)(x,(c,d))
  &=2A+B+(2C+D)\ind_{\{b=d\}}.\label{eq:sum-first}
\end{align}
For an arbitrary matrix $M$, \eqref{eq:decomp} gives
\[
(MP_N)(x,y)=sM(x,y)
+c_N\sum_{v\in\Z/N\Z}M(x,(c,v))
+c_N\sum_{u\in\Z/2\Z}M(x,(u,d)).
\]
Using \eqref{eq:sum-second}--\eqref{eq:sum-first}, and
$s+c_NN+2c_N=1$, $s+c_NN=q$, $s+2c_N=r$, we obtain
\begin{equation}\label{eq:recurrence}
F(A,B,C,D)P_N
=F\bigl(A+c_N(B+C),\ qB+c_ND,\ rC+c_ND,\ sD\bigr).
\end{equation}

Define coefficients, for all integers $t\ge0$, by
\begin{align*}
A_t&=\frac{1-q^t-r^t+s^t}{2N},&
B_t&=\frac{q^t-s^t}{N},\\
C_t&=\frac{r^t-s^t}{2},&
D_t&=s^t.
\end{align*}
At time zero these are $(0,0,0,1)$, so $F(A_0,B_0,C_0,D_0)=I$.
They satisfy precisely the four recurrences
\[
\begin{split}
A_{t+1}&=A_t+c_N(B_t+C_t),\qquad
B_{t+1}=qB_t+c_ND_t,\\
C_{t+1}&=rC_t+c_ND_t,\qquad
D_{t+1}=sD_t.
\end{split}
\]
For example, $q-1/2=s$ gives the recurrence for $B_t$, and
$r-1/N=s$ gives that for $C_t$. For the constant coefficient,
\[
A_{t+1}-A_t
=\frac{(1-q)q^t+(1-r)r^t-(1-s)s^t}{2N}
=c_N\left(\frac{q^t-s^t}{N}+\frac{r^t-s^t}{2}\right),
\]
using $1-q=1/N$, $1-r=1/2$, and $1-s=1/2+1/N$.
The recurrence for $D_t$ is immediate.
Induction with \eqref{eq:recurrence} therefore proves
$P_N^t=F(A_t,B_t,C_t,D_t)$ for every $t\ge0$.

Rearranging this identity using $\sigma=2U-1$ proves the promised formula:
\begin{proposition}\label{prop:powers}
For every $t\in\N$ and every pair of states $x=(a,b)$, $y=(c,d)$,
\begin{equation}\label{eq:powers}
P_N^t(x,y)
=\frac{1+\sigma q^t+(NV-1)(r^t+\sigma s^t)}{2N}.
\end{equation}
In particular this formula holds at $t=0$ and equals the identity kernel
there.
\end{proposition}

One useful marginal follows either from the coefficient formula or by
summing \eqref{eq:powers} over $d$:
\begin{equation}\label{eq:blockmass}
\sum_{d\in\Z/N\Z}P_N^t(x,(c,d))
=\frac{1+\sigma q^t}{2},
\end{equation}
because $\sum_d(N\ind_{\{b=d\}}-1)=0$.

\section{Exact total variation, including time zero}

The signs inside the absolute values are essential. We establish them
before summing; no norm approximation or spectral bound replaces the
actual TV distance.

\begin{lemma}\label{lem:sign-estimates}
For each integer $t\ge1$,
\[
r^t+s^t\le q^t,\qquad (N-1)(r^t-s^t)\le q^t.
\]
\end{lemma}
\begin{proof}
We have $q=r+s$ and $r,s\ge0$. The first inequality follows by expanding
$(r+s)^t$ and retaining its two endpoint terms. For the second, use
$(N-1)(r-s)=q$ and the difference-of-powers identity:
\[
(N-1)(r^t-s^t)
=q\sum_{j=0}^{t-1}r^{t-1-j}s^j
\le q(r+s)^{t-1}=q^t.
\]
The middle inequality holds because every binomial coefficient in
$(r+s)^{t-1}$ is at least one and all monomials are nonnegative. This also
covers $t=1$, for which the sum has its single term equal to one. The Lean
proof establishes the same inequalities by elementary induction.
\end{proof}

Fix a starting state $x=(a,b)$ and $t\ge1$. Subtracting $1/(2N)$ from
\eqref{eq:powers} gives the numerator
\[
\sigma q^t+(NV-1)(r^t+\sigma s^t).
\]
In the block $c=a$ this is nonnegative: for $V=1$ both summands are
nonnegative, while for $V=0$ it is $q^t-r^t-s^t\ge0$.
In the other block it is nonpositive: for $V=1$ it is
$-q^t+(N-1)(r^t-s^t)\le0$, while for $V=0$ it is
$-q^t-r^t+s^t\le0$, since $s\le r$.
These signs remain valid after division by the positive number $2N$.

By \eqref{eq:blockmass}, the sum of the stationary deviations in the
starting block is
\[
\frac{1+q^t}{2}-\frac12=\frac{q^t}{2}.
\]
In the other block it is $-q^t/2$. The absolute deviations therefore sum
to $q^t$, and
\begin{equation}\label{eq:tv-positive}
\frac12\sum_y\left|P_N^t(x,y)-\pi_N(y)\right|=\frac{q^t}{2}
\qquad(t\ge1).
\end{equation}
This equality holds for every starting state $x$, so it also gives their
maximum. At time zero the row of $P_N^0$ is a point mass, and direct summation
instead gives
\[
\frac12\left(1-\frac1{2N}+(2N-1)\frac1{2N}\right)
=1-\frac1{2N}.
\]
Together these prove \eqref{eq:tv-main}.

\section{True mixing times and the absence of cutoff}

\subsection{Existence and least-time properties}

For fixed $N$, $0<q<1$, so $q^t\to0$. For completeness, this follows
without logarithmic asymptotics: set $a=q^{-1}-1>0$. Bernoulli's inequality
gives $(1+a)^t\ge1+ta$, and hence
\[
0<q^t\le\frac1{1+ta}\longrightarrow0.
\]
By \eqref{eq:tv-positive}, for every $\varepsilon>0$ some positive integer
$t$ satisfies $d_N(t)\le\varepsilon$. The threshold set is nonempty and,
by well-ordering of $\N$, has a minimum. In particular,
\begin{equation}\label{eq:minimum}
d_N(t_N(\varepsilon))\le\varepsilon,
\qquad u<t_N(\varepsilon)\ \Longrightarrow\ d_N(u)>\varepsilon.
\end{equation}
These are properties of the walk's actual least mixing times, not an
assumed or substituted scalar sequence.

\subsection{Two exact thresholds and a divergent mixing time}

Since $N\ge4$,
\[
d_N(0)=1-\frac1{2N}>\frac34,
\qquad d_N(1)=\frac q2<\frac12<\frac34.
\]
Thus $t_N(3/4)=1$ exactly. Likewise $T_N=t_N(1/4)$ is positive, so
attainment in \eqref{eq:minimum} and \eqref{eq:tv-positive} imply
$q^{T_N}\le1/2$. Bernoulli's inequality, now with increment $-1/N$, gives
\[
1-\frac{T_N}{N}\le\left(1-\frac1N\right)^{T_N}
=q^{T_N}\le\frac12.
\]
It follows that
\begin{equation}\label{eq:T-lower}
T_N\ge\frac N2\ge2.
\end{equation}
In particular $T_N\to\infty$, while the high-tolerance mixing time stays
equal to one.

\subsection{The actual transition width and quantified cutoff}

Our width is the difference between the genuine threshold times:
\[
W_N=T_N-t_N(3/4)=T_N-1\ge1.
\]
Consequently
\begin{equation}\label{eq:width-ratio}
\frac{T_N}{W_N}=\frac{T_N}{T_N-1}
=1+\frac1{T_N-1}\longrightarrow1.
\end{equation}
For every $N\ge4$ this ratio is greater than one and at most two. In
particular, it cannot tend to infinity. The time needed to move from TV
distance $3/4$ to $1/4$ has the same order as the ordinary mixing time;
it is not a window of vanishing relative width.

The quantified standard cutoff test fails independently at the single
allowed tolerance $\varepsilon=1/4$:
\[
\frac{t_N(1/4)}{t_N(1-1/4)}
=\frac{T_N}{1}=T_N\longrightarrow\infty,
\]
whereas \eqref{eq:cutoff} would require this ratio to tend to one. This
proves all remaining assertions of Theorem~\ref{thm:main}.

\begin{corollary}\label{cor:refutation}
Under the ordinary spectral-gap and standard total-variation meanings,
divergence of the proposed multiplicity-to-gap ratio is not sufficient for
cutoff on lazy abelian-group random walks. This holds whether ``minimal
nontrivial eigenvalue'' means the smallest nontrivial eigenvalue of $P_N$
or the smallest positive eigenvalue of $I-P_N$.
\end{corollary}
\begin{proof}
A universal sufficient condition would apply to every family in the
asserted class, hence to \eqref{eq:kernel}. The support, group, and stochastic
properties above prove that this is such a family, even with exact half
laziness and uniform supported steps. Equation~\eqref{eq:ratios} proves
both candidate antecedents. Equations~\eqref{eq:width-ratio} and
\eqref{eq:cutoff}'s failure prove the corresponding conclusions false.
Thus this one fully specified witness refutes the universal sufficient
direction. Refuting that direction is enough to refute the proposed iff
under these meanings; the necessary direction need not be decided.
\end{proof}

\section{Correspondence with the frozen Lean certificate}

The formalization uses the same index $N=n+4$, state group, support,
transition kernel, and threshold conventions as this report. Its principal
definitions are summarized below. Every spectral and mixing formula used
as a conclusion in the preceding proof is proved from these definitions.

\begin{center}\small
\begin{tabularx}{\textwidth}{@{}p{0.30\textwidth}X@{}}
\toprule
Lean definition&Mathematical object\\
\midrule
\code{State}, \code{support}&$G_N$ and the actual generating set $S_N$.\\
\code{jumpLaw}, \code{transition}&$\mu_N$ and $P_N(x,y)=\mu_N(y-x)$.\\
\code{tvFrom}, \code{worstTV}&The absolute-value finite TV sum from each start and its finite maximum.\\
\code{Characters}&Actual complex additive characters of $G_N$.\\
\code{fourierEigenvalue}&The Fourier sum of the actual increment law, proved to be its matrix eigenvalue.\\
\code{complexTransition}&Entrywise complexification of the already defined real matrix $P_N$.\\
\code{spectralGap}&One minus the finite maximum nontrivial character eigenvalue. All values are proved real.\\
\code{minTransitionEigenvalue}&The finite minimum of the actual nontrivial transition eigenvalues.\\
\code{minTransitionMultiplicity}&Root multiplicity in the actual transition characteristic polynomial at that minimum.\\
\code{complexComplement}&The actual matrix $I-P_N$.\\
\code{gapModeMultiplicity}&Root multiplicity in the actual complement characteristic polynomial at the proved ordinary gap.\\
\code{mixingTime}&The natural-number infimum of the actual TV threshold set, with nonemptiness, attainment, and minimality proved.\\
\code{totalTime}, \code{window}&Real-valued $T_N$ and $T_N-t_N(3/4)$.\\
\code{HasTVCutoff}&The quantified condition \eqref{eq:cutoff}.\\
\bottomrule
\end{tabularx}
\end{center}

The proof modules divide the obligations as follows.
\begin{itemize}
\item \module{Basic.lean} gives the actual objects, normalizations, and
elementary inequalities. \module{Walk.lean} proves support generation,
the normalized uniform Cayley formula, stochasticity, stationary mass,
stochastic powers, exact holding, and positive two-step probabilities.
\item \module{Spectrum.lean} constructs the equivalence between factor
characters and all group characters. It proves actual matrix action and
diagonalization in \code{AddChar.complexBasis}, both characteristic
polynomials, complete root classifications, finite extrema, and the
genuine algebraic multiplicities.
\item \module{Powers.lean} proves the matrix recurrence and actual power
formula, the sign inequalities, and the finite absolute-value TV sums,
including time zero and the maximum over all starting states.
\item \module{MixingBase.lean} defines the actual least-threshold times
and proves their properties and limits. Its intermediate lemmas explicitly
take the exact TV profile as a parameter. In the final theorem every such
parameter is discharged with \code{worstTV_eq}, the identity proved in
\module{Powers.lean}; it is not an assumption of the final result.
\item \module{Counterexample.lean} assembles the closed certificate and
four sufficient-condition refutations. \module{Audit.lean} requests the
axiom lists of the five capstone theorems. The top-level library imports the
counterexample module.
\end{itemize}

All capstones lie in the namespace \code{Conjecture1421}:
\begin{enumerate}[label=\arabic*.]
\item \code{verified_counterexample} proves
\code{CounterexampleCertificate}, a proposition exposing the actual walk,
group-order divergence, support, probability, character, spectral,
matrix-power, TV, minimum-time, window, and cutoff obligations together.
It has no mathematical premise.
\item \code{literal_sufficient_criterion_false} negates the implication
from the minimum-transition multiplicity ratio's divergence to
\code{HasTVCutoff} for this actual family.
\item \code{gap_mode_sufficient_criterion_false} does the same for the
minimum positive eigenvalue of the actual $I-P_N$.
\item \code{literal_window_sufficient_criterion_false} negates the
implication from the first ratio's divergence to divergence of the actual
time/window ratio.
\item \code{gap_mode_window_sufficient_criterion_false} does the same
for the second ratio.
\end{enumerate}

The four corollaries are deliberately statements about this fully
specified witness; they do not pretend to encode every undefined phrase
of the original or an entire class of arbitrary random walks. The closed
certificate establishes membership in the proposed class and the proved
antecedents. Corollary~\ref{cor:refutation} explains the logical step from
this witness to failure of universal sufficiency. No profile assumption,
spectral formula, multiplicity, or mixing-time bound remains as a hypothesis
of any capstone.

\section{Reproduction and supplemental verification}

\subsection{Pinned formal project}

The project requests Lean \code{leanprover/lean4:v4.19.0} and Mathlib
\code{v4.19.0}. Its dependency manifest pins Mathlib to
\begin{center}\small\ttfamily
c44e0c8ee63ca166450922a373c7409c5d26b00b.
\end{center}
The manifest also pins all transitive dependencies. Warnings are treated
as errors in the project configuration. From the directory containing
\module{lakefile.toml}, with the pinned toolchain and dependencies
available, the reproduction commands are
\begin{quote}\ttfamily\small
lake build\\
lake env lean -DwarningAsError=true Conjecture1421/Audit.lean
\end{quote}
The audit command imports the actual capstones rather than proving an
unrelated sample theorem. The author-stage fresh direct Lean replay of all
eight modules succeeded with warnings treated as errors. Its capstone
axiom output consists only of the standard Lean logical axioms
\code{propext}, \code{Classical.choice}, and \code{Quot.sound}.
The source has no admitted proof, added axiom, native-decision proof, or
custom unsafe proof dependency.

The separate isolated engineering run completed a source-based Lake build
and eight strict module replays successfully. Its audit counted 342 owned
declarations, 292 safe owned declarations, 31,861 declarations in the
combined dependency closure, and 3,701 imported modules. The five capstones
again had only the three standard axioms listed above. An independent
finite-regression rerun reproduced the recorded output exactly. A positive
control, five negative controls, and seven build-harness hardening fixtures
also completed as expected. The independent engineering review passed after
explicitly classifying all 50 generated unsafe runtime helpers (42 definitions
and eight compiler placeholder axioms). None occurs in any safe owned proof
closure or in any of the five capstones. The raw driver status, which leaves
this classification pending, is preserved alongside the completed separate
classification and verdict. All 62 command receipts are retained, including
six intentional nonzero negative-control results; all stderr streams are empty.

These checks are local verification, not maintainer acceptance. Cached generic
library artifacts and pinned source bytes were identified separately; this is
not a full Mathlib rebuild or proof of source-to-cache provenance. Lean was
not bootstrapped. The dependency audit and graph consistency checks are not a
second proof kernel; the fresh source compilations supply kernel checking
relative to the supplied runtime and imports. PATH network wrappers are not
operating-system isolation. Complete source and configuration review complements
these checks. Detailed identities, controls, and limits are preserved in the
accompanying verification records. The finite evidence below remains
supplementary to the mathematical proof.

\subsection{Exact finite regression}

The supplemental \module{exact_regression.py} uses exact rational
arithmetic. It independently constructs increments and transition
probabilities and evolves distributions by direct multiplication. Its
recorded run covered
\[
N\in\{4,5,6,7,8,11,16\},\qquad t=0,1,\ldots,24,
\]
all starting states, and $56{,}700$ actual kernel entries. It checked the
power identities, actual absolute-value TV distances, least threshold
times, and an explicit real eigenbasis with its rank and eigenvalue action.
The run also included three intentional negative controls. To reproduce
these supplemental checks, run \code{python3 exact_regression.py} from the
directory containing the script.

These finite checks are corroborating arithmetic evidence only. They do
not establish a universal-in-$N$ or universal-in-$t$ identity, an asymptotic
limit, or a Lean theorem. The proofs in this report and the closed formal
certificate supply those claims without using the regression output as
an assumption.

\section{Conclusion and scope of the refutation}

The actual family \eqref{eq:kernel} has growing finite abelian state groups,
uniform symmetric supported steps, exact half laziness, and diverging
ordinary mixing times. Its full spectra show that both proposed
multiplicity-to-ordinary-gap ratios tend to infinity. Its actual
total-variation profile and least threshold times show simultaneously that
standard TV cutoff fails and the genuine time/window ratio tends to one.
This refutes the explicit sufficient direction, and therefore the explicit
iff, under the stated ordinary-gap and total-variation meanings.

The interpretation limits remain part of the result. The source does not
specify a mixing metric or an exact window convention, and it does not
define a conditional spectral gap, second-order behavior, or a particular
Diaconis--Saloff-Coste comparison characterization. We have not formalized
those undefined phrases or claimed a result about every possible mixing
metric or an unstated additional hypothesis. Both natural meanings of
``minimal nontrivial eigenvalue'' have, however, been handled by actual
operator spectra and actual algebraic multiplicities, with no conflation
between the two.

\begin{thebibliography}{1}
\bibitem{levin}
David A. Levin, \emph{Cutoff for Markov Chains}, LMS--EPSRC Durham Symposium,
2017, slide~23 (PDF page~23).
\href{https://pages.uoregon.edu/dlevin/TALKS/durham.pdf}{University of Oregon slide deck}.
Cited only for the standard cutoff convention, after the independent
derivation of this counterexample.
\end{thebibliography}

\end{document}
